\chapter{Execution Map}

This chapter follows only the committed default benchmark
\code{benchmarks/3d\_hetero\_ssr\_default/case.toml}. In that benchmark the
main path is:
\[
\text{3D hetero SSR} \rightarrow
\text{indirect continuation} \rightarrow
\text{owned-row assembly} \rightarrow
\text{PMG shell backend}.
\]

\section{Benchmark Definition}

For the current mainline reading path, the important settings are:
\begin{itemize}
\item \code{analysis = ssr},
\item \code{method = indirect},
\item \code{elem\_type = P4},
\item \code{seepage = false},
\item \code{node\_ordering = block\_metis},
\item \code{mpi\_distribute\_by\_nodes = true},
\item \code{constitutive\_mode = overlap},
\item \code{tangent\_kernel = rows} by execution default,
\item \code{solver\_type = PETSC\_MATLAB\_DFGMRES\_HYPRE\_NULLSPACE},
\item \code{pc\_backend = pmg\_shell}.
\end{itemize}

Everything else in this document is organized around that exact path. Generic
framework defaults and other runner surfaces are moved to
Appendix~\ref{app:execution}.

\section{MATLAB Driver to PETSc Runner}

The MATLAB reference script
\matlabpkg{slope\_stability\_3D\_hetero\_SSR.m} concentrated the whole solve in
one file. The current PETSc path is the same algorithm broken into explicit
stages:
\begin{enumerate}
\item \petscmod{core/run\_config.py} loads the TOML case and resolves paths.
\item \petscmod{cli/run\_case\_from\_config.py} selects the 3D SSR runner.
\item \petscmod{cli/run\_3D\_hetero\_SSR\_capture.py::run\_capture} performs
mesh preprocessing, constitutive-builder setup, solver construction, and
continuation launch.
\item \petscmod{continuation/indirect.py::SSR\_indirect\_continuation} drives
the branch tracking.
\item \petscmod{nonlinear/newton.py::newton\_ind\_ssr} performs the nested
Newton solves used at each continuation step.
\item \petscmod{linear/solver.py::SolverFactory.create} provides the PMG-shell
preconditioned outer Krylov solve.
\end{enumerate}

\section{Mainline Call Chain}

For the default 3D PMG benchmark, the runner sequence is:
\begin{enumerate}
\item load the mesh with \code{load\_mesh\_from\_file(...)},
\item reorder nodes with \code{reorder\_mesh\_nodes(...)} using
\code{block\_metis},
\item expand material rows with \code{heterogenous\_materials(...)},
\item choose the owned distributed assembly path through
\petscmod{cli/assembly\_policy.py},
\item build owned elastic rows with
\code{assemble\_owned\_elastic\_rows\_for\_comm(...)},
\item create \code{ConstitutiveOperator},
\item prepare the owned tangent pattern with
\code{prepare\_owned\_tangent\_pattern(...)},
\item create the linear solver through \code{SolverFactory.create(...)},
\item launch \code{SSR\_indirect\_continuation(...)}.
\end{enumerate}

That list is the modern replacement for reading the MATLAB script from top to
bottom.

\section{What Moved Out of the MATLAB Driver}

The main structural change is not a new nonlinear algorithm. It is that the old
driver responsibilities are now explicit modules:

\begin{longtable}{@{}p{0.30\textwidth}p{0.30\textwidth}p{0.34\textwidth}@{}}
\toprule
MATLAB role & MATLAB location & Mainline PETSc location \\
\midrule
\endhead
Case declaration & top of \matlabpkg{slope\_stability\_3D\_hetero\_SSR.m} & \code{case.toml} plus \petscmod{core/run\_config.py} \\
Mesh and material preprocessing & script body plus \matlabpkg{+MESH/*} and \matlabpkg{+ASSEMBLY/*} & \petscmod{cli/run\_3D\_hetero\_SSR\_capture.py}, \petscmod{mesh/*}, \petscmod{fem/*} \\
Constitutive handle construction & \matlabpkg{+CONSTITUTIVE\_PROBLEM/CONSTITUTIVE.m} & \petscmod{constitutive/problem.py::ConstitutiveOperator} \\
Linear solver setup & \matlabpkg{+LINEAR\_SOLVERS/set\_linear\_solver.m} & \petscmod{linear/solver.py::SolverFactory.create} \\
Continuation launch & \matlabpkg{+CONTINUATION/SSR\_indirect\_continuation.m} & \petscmod{continuation/indirect.py::SSR\_indirect\_continuation} \\
\bottomrule
\end{longtable}

\section{Shortest Reading Order}

For the current mainline path, the shortest route through the PETSc code is:
\begin{enumerate}
\item \petscmod{cli/run\_3D\_hetero\_SSR\_capture.py} for orchestration,
\item \petscmod{cli/assembly\_policy.py} for the owned-path choice,
\item \petscmod{fem/distributed\_elastic.py} and
\petscmod{fem/distributed\_tangent.py} for matrix construction,
\item \petscmod{constitutive/problem.py} for the constitutive builder,
\item \petscmod{linear/solver.py} for the PMG-shell solver path,
\item \petscmod{nonlinear/newton.py} and
\petscmod{continuation/indirect.py} for branch advancement.
\end{enumerate}

\section{Where To Look For Inactive Paths}

This chapter intentionally excludes generic framework defaults, seepage
runners, and other entry points that are not used by the default parallel PMG
benchmark. Those are collected in Appendix~\ref{app:execution}.
